\section{Discussion}\label{sec:discussion}

\subsection{What the range means for preparers, auditors and validators}

\paragraph{Preparers.} SAB 119 expects a registrant to document an identified
range of expected credit losses and its reasons for the point it books. The
results suggest the range should include the model. On this book a median
of \num{res.rel.median}\% of the allowance separates models within one standard
error of the best held-out log loss (\num{ph.paired2.rel}\% under a strict
paired test), and retraining the chosen model on another random sample of the
same size moves it at least as far. A lender can measure both with the procedure
here: fit a small grid, keep the models within a stated tolerance of the best
held-out log loss, refit the best on independent samples of its data, and
report the allowance across both. Bootstrap refits are a poor substitute:
they are less accurate than the model they refit and biased upward. None of this needs data the lender does not already have.

\paragraph{Auditors.} AS 2501 asks the auditor to evaluate the company's
selection among reasonable alternatives, and AS 2810 treats a recorded
estimate outside a reasonable range as misstated by the distance to the closest
reasonable estimate. Equally accurate models are reasonable alternatives on
any accuracy-based standard. A set like the one used here gives the auditor a
reproducible range to set the booked estimate against, in place of a judgment
about whether a single model is ``reasonable''.

\paragraph{Model validators.} The 2026 interagency guidance asks validators to
benchmark against alternative models, but it does not address accounting
estimates. Benchmarking that compares a chosen model with alternatives only on
accuracy will pass every model in the set while missing the
\num{res.rel.median}\% spread in what the models are for. Validation of an
allowance model should compare allowances, not only accuracy, across the
models the bank could defensibly have chosen, and include refit noise.

\paragraph{Monotone constraints are not a fix.} Constraining the direction of
credit risk narrowed the range at barely more than half the dates (H4), and
not at all on average. The lifetime projection runs through inputs the
constraints leave free.

\paragraph{No range is forecast uncertainty.} At all but one of the
dates with eight or more years of outcomes, the realized loss on the scored
loans lay outside the set's range (H5), and outside the range of the refits at
all but \num{p2.in_refit} (bootstrap) or \num{p2.in_rsm} (independent samples). Before the crisis every model in the
set, fed a forecast that held house prices rising at their 2005 or 2006 pace
and a severity table from benign years, missed the losses by a wide margin;
after it the allowance overshot. Before the crisis the miss lay in the
forecast and the severity table rather than in the default model, whose
lifetime PDs on the realized path were close to the realized default rates
(\cref{tab:decomp}), and the models agreed with one another far more than
with the outcome. A range across models documents one kind of uncertainty
and should not be read as a confidence interval for losses.

\subsection{Limitations}\label{sec:limits}

\begin{itemize}[leftmargin=*, itemsep=2pt]
\item \emph{One book, one product.} The data are Freddie Mac's standard dataset
of conforming, fully amortizing fixed-rate mortgages. Banks' books contain
other products, and an agency guarantor bears credit risk differently from a
bank holding a loan. The magnitudes should not be read as a bank's; the
procedure carries over.
\item \emph{No delinquency input.} The design has no current-delinquency
input. Loans 30 to 179 days past due are \num{dq.share.min}\% to
\num{dq.share.max}\% of the portfolio balance (most in December
\num{dq.share.max_year}), and they default far more often than current loans:
in 2010, \num{dq.2009.d90}\% of the loans 90 to 179 days past due at December
2009 defaulted and \num{dq.2009.d30}\% of those 30 to 89 days past due, against
\num{dq.2009.cur}\% of current loans. Any lender would use delinquency, which
would raise the level of the allowance, most in 2008 to 2011. Whether it
would narrow the range across models is not tested here.
\item \emph{The tolerance.} The recorded set admits models within one
bootstrap standard error of the best model's own log loss, which is generous:
a test on paired differences admits a median of \num{ph.paired2.size} models at
two standard errors. The relative range is still \num{ph.paired2.rel}\% of the
allowance under that strict set.
\item \emph{The model grid.} The set contains only the families tried here.
A wider grid (neural networks, other link functions) could widen the range,
which is therefore a lower bound on model-choice uncertainty within these
families.
\item \emph{The policy range.} At annual resolution two of the three reversion
policies coincide, and the length of the forecast period is not varied, so
H2 compares the model range with a narrow slice of the policy discretion the
standard allows.
\item \emph{Data as revised.} Each date uses only data dated on or before $T$,
but in its 2026 form: (a) Freddie Mac's Release 47 performance history and actual losses,
which may correct what was reported at the time and include recoveries and
expenses posted after a loan's zero-balance date, so the severity table at $T$
contains amounts a lender at $T$ had not yet seen; (b) current FRED vintages
of the house-price index and unemployment rates, which are revised; (c)
long-run means that include year $T$ itself; (d) the exclusion from the
severity table of loans repurchased for a defect found after $T$. None of
these is a default outcome after $T$, but (a), (b) and (d) are information a
lender at $T$ did not have in this form.
\item \emph{Censored loan-years.} Loan-years ending in removal for a defect or
a reperforming sale without a prior default are dropped from training, which
treats that removal as uninformative about default.
\item \emph{Pandemic forbearance.} By the recorded definition, a loan 180 days
past due in forbearance has defaulted. Models trained after 2020 learn that
spike, and most of those loans were later paid off or remain active. This
raises lifetime PDs at the later dates; the severity table offsets it in part
because payoffs after default count as zero-loss resolutions.
\item \emph{Severity.} The loss-given-default table is held fixed across
models, uses only defaults resolved by the reporting date, and excludes
modification costs, which Freddie Mac's actual-loss field omits. The design's
secondary analysis that lets severity vary with the default model was not run
(\cref{app:deviations}); a model for severity would add another source of
variation.
\item \emph{Realized losses.} Outcomes end in June 2026, so the realized loss
of the later portfolios is truncated. Defaults still unresolved then are
valued at the severity of resolved defaults, which is most of realized loss
from 2015 (\cref{tab:decomp}); the lower bound counts them at zero. Neither
changes the verdict on H5.
\end{itemize}

\subsection{Conclusion}

The CECL standard leaves the method to the lender, and the evidence here is
that the choice is not innocuous. Equally accurate default models trained on
the same history give allowances a median \num{res.rel.median}\% apart, more
than the gap between one- and three-year reversion, and no more than
retraining one model on another sample of its loans. Neither source of
variation tells a lender where losses will land. Documenting the range from
the model, as SAB 119 already asks for the estimate, is cheap with the
procedure released here, and it is the part of the identified range a bank
can measure with its own data.
